\documentclass{res} 
\usepackage{microtype}
\usepackage{newpxmath}
\usepackage{url}
\renewcommand{\rmdefault}{ppl}

\begin{document}
%\vspace*{-12pt}
%\moveleft.5\hoffset\centerline{\large\bf Devlin Mallory} % Your name at the top
%\vskip-3pt
%%\moveleft\hoffset\vbox{\hrule width 6.90in height .8pt} % Horizontal line after
\newdimen\linew
\linew=\textwidth
\advance\linew\hoffset 
%\fulline
%\moveleft.5\hoffset\centerline{}

\name{Devlin Mallory\\\vspace{-10pt}}
\address{dmallory@bcamath.org\qquad \url{https://devlin-mallory.github.io/} 
%\qquad 831.234.5488 
\\\vspace{-2pt} }
\setlength{\leftmargini}{-10pt}

\renewcommand\sectionfont{\large \sc}
\begin{resume}
{\bf Interests:} Singularities, Frobenius summands, differential operators, positivity of vector bundles, jet~schemes, Macaulay2.

\section{Employment} 
\ \\
{\bf Basque Center for Applied Mathematics}, Marie Sk\l odowska-Curie Fellow\hfill June 2026 -- Present
\\\\
{\bf Basque Center for Applied Mathematics}, Postdoctoral Researcher\hfill October 2024 -- May 2026
\\\\
{\bf University of Utah}, Wylie Assistant Professor (Lecturer)\hfill Fall 2021 -- Summer 2024
\\\\
{\bf University of Michigan}, NSF Graduate Fellow\hfill Fall 2015 -- Spring 2021
\section{Education} 
\ \\
{\bf University of Michigan}\hfill Fall 2015 -- Spring 2021
\\ Ph.D., Advisor: Mircea Musta\c{t}\u{a}
\\ Thesis: \emph{Singularities of Birational Geometry via Arcs and Differential Operators}
\\
\\
{\bf University of California, Berkeley} \hfill Fall 2009 -- Fall 2013\\B.A. Pure Mathematics (Honors),  B.A. Physics 
 %{\sl Summa cum laude} (GPA: 3.97) %\hfill Fall 2009 - Fall 2013 % \\ Pure Mathematics (Honors), Physics  \\



\section{Papers and preprints} 
%\bibliography{devlin}
\begin{enumerate}
%\item Feliks? Undergrads? Suchitra?
\item Dg models for derived differential operators (with Jack Jeffries, Claudia Miller, Josh Pollitz, and Eamon Quinlan-Gallego). arXiv:2610.23778, 2026.
\item The WittVectors package for Macaulay2 (with Anne Fayolle, Abhay Goel, Eamon Quinlan-Gallego, and Teppei Takamatsu). arXiv:2609.37931, 2026.
\item Computing direct sum decompositions (with Mahrud Sayrafi). Journal of Symbolic Computation 133, 2026.
\item Antiampleness and ampleness of the Frobenius cokernel. arXiv:2510.03193, 2025.
\item The tilting property for $F_*^e\mathcal O_X$ on Fano surfaces and threefolds. arXiv:2405.14070, 2024.
\item Homogeneous coordinate rings as direct summands of regular rings. Illinois Journal of Mathematics 68(1), 2024.
\item Finite $F$-representation type for homogeneous coordinate rings of non-Fano varieties. \'Epijournal de G\'eom\'etrie Alg\'ebrique 7, 2023.
\item An explicit self-duality (with Nikolas Kuhn, Vaidehee Thatte, and Kirsten Wickelgren). In \emph{Stacks Project Expository Collection}, LMS Lecture Note Series, number 480, Cambridge University Press, 2023.
\item Bigness of the tangent bundle of del Pezzo surfaces and D-simplicity. Algebra \& Number Theory 15(8), 2021.
\item
Minimal log discrepancies of determinantal varieties via jet schemes. Journal of Pure and Applied Algebra 225(2), 2021.
\item
Triviality of arc closures and the local isomorphism problem. Journal of Algebra 544, 47--61, 2020.
\end{enumerate}

\section{Invited talks}
\begin{enumerate}
\item 
%``F-decomposition matrices for cyclic quotient singularities''
Seminario de \'Algebra Conmutativa y Geometr\'{\i}a Algebraica, Centro de Investigaci\'on en Matem\'aticas, Guanajuato, November 2026.
\item
%``F-signatures and Frobenius alpha invariants of del Pezzo surfaces''
59th Congreso Nacional de la Sociedad Matem\'atica Mexicana, Special Session on Commutative Algebra and its Applications, Tlaxcala, October 2026.
\item
%``Studying singularities via the complexity of the Frobenius''
University of the Basque Country (UPV/EHU), Algebra Seminar, December 2025.
\item
%``Positivity of tangent sheaves and the study of Frobenius pushforwards''
Universit\`a di Genova, Algebra \& Geometry Seminar, November 2025.
\item
%``Computing direct sum decompositions (and some applications)''
UCLA Birational Geometry Seminar, June 2025.
\item 
%NEEDS TITLE
Meeting of the Spanish Network for Algebraic Geometry and Singularity Theory, Madrid, January 2025.
\item
%``Finite F-representation type for homogeneous coordinate rings''
Commutative Algebra and Algebraic Geometry in Tucson, May 2024.
\item 
%NEEDS TITLE
Stony Brook University, Algebraic Geometry Seminar, April 2024.
\item
%``Understanding singularities via the complexity of the Frobenius''
MATH For All Conference, Salt Lake City, April 2024.
\item 
%NEEDS TITLE
MSRI/SLMath, Special Seminar on Finite $F$-Representation Type for Rings and Varieties, March 2024.
\item
%``Finite F-representation type for homogeneous coordinate rings''
Joint Mathematics Meetings, Special Session on Recent Developments in Commutative Algebra, January 2024.
\item 
%``Finite F-representation type for homogeneous coordinate rings''
AMS Fall Central Sectional, Special Session on Commutative Algebra, Differential Operators, and Singularities, October 2023.
\item 
%``Direct Summands of Regular Rings''
University of Minnesota Commutative Algebra and Algebraic Geometry Seminar, April 2022.
\item 
%``Differential operators on singular rings and the global geometry of projective varieties,''
University of Connecticut Algebra Seminar, March 2021.
\item 
%``Bigness of the tangent bundle and differential properties of singularities,''
University of Utah Algebraic Geometry Seminar, October 2020.
\item 
%``D-simplicity in characteristic 0 and characteristic p,''
AMS Fall Eastern Sectional, Special Session on Algebraic Singularities in Arbitrary Characteristic, October 2020.
\item 
%``D-simplicity in characteristic 0 and characteristic p,''
Differential Operators in Commutative Algebra Seminar, May 29 and June 5, 2020.
\item 
%``Triviality of jet closures and the local isomorphism problem,''
University of Michigan Commutative Algebra Seminar, October 2019.
\item 
%``Minimal log discrepancies of determinantal varieties via jet schemes,''
AMS Fall Central Sectional, AMS Contributed Paper Session on Algebra, Geometry and Lie Theory, September 2019.
\item (poster presentation) 
%``Arc closures and the local isomorphism problem,''
Midwest Algebraic Geometry Graduate Conference, March 2019.
\end{enumerate}



\section{Organizational activities}
\begin{enumerate}
\item Teaching Assistant, RGAS Summer School on Perfectoid Techniques, Summer School of the Red de Geometr\'{\i}a Algebraica y Singularidades, May 2025.
\item Teaching Assistant, Methods in Mixed Characteristic Geometry, Fall School at JGU Mainz, October 2024.
\item Organizer, Macaulay2, Computational Algebraic Geometry and String Theory workshop and conference, held at the University of Utah, Spring 2024.
\item Organizer, Utah Algebraic Geometry Seminar, Fall 2022 -- Spring 2024.
\item Organizer, Introduction to Algebraic Geometry Reading Group, Winter 2023.
\item Teaching Assistant, Master Class: Local Cohomology and related topics, Stockholm, August 2022.
\item Mentor, University of Michigan Directed Reading Program, Summer 2020.
\item Organizer, Student Algebraic Geometry Seminar, Fall 2018 -- Winter 2021.
\item Organizer, Seminar on singularities in algebraic geometry and commutative algebra, Fall 2018 -- Winter 2019.
\item Organizer, Summer graduate student minicourse series, 2018 -- 2021.
\end{enumerate}

 
\section{Mathematical software} 

\begin{enumerate}
\item \texttt{WittVectors} (with Anne Fayolle, Abhay Goel, Eamon Quinlan-Gallego, and Teppei Takamatsu). Macaulay2 package for computations with Witt vectors, Frobenius lifts, and quasi-$F$-splittings. Distributed with Macaulay2 since version 1.26.05.
\item \texttt{DirectSummands} (with Mahrud Sayrafi). Macaulay2 package for decompositions of modules and coherent sheaves. Distributed with Macaulay2 since version 1.26.06.
\item \texttt{Varieties} (with John Cobb, Ritvik Ramkumar, Mahrud Sayrafi, Gregory G.~Smith, and Keller VandeBogert). Macaulay2 package for working with affine and projective varieties and coherent sheaves on them. Distributed with Macaulay2 since version 1.23.
\end{enumerate}

\section{Other professional activities}
\begin{enumerate}
\item Invited participant, Collaborate@ICERM, January 2025.
\item Invited participant, IAS Summer Collaborators Program, July 2024.
\item Reviewer, AMS Mathematical Reviews, 2023 -- Present.
\item Reviewer, zbMATH Open, 2022 -- Present.
\end{enumerate}

\section{Awards}
\begin{enumerate}
\item Marie Sk\l odowska-Curie Postdoctoral Fellowship, June 2026. Project acronym: SIPOCAG.
\item University of Utah Outstanding Postdoc Award, April 2023.
\item University of Michigan Math Department Graduate Student Award, for ``contributions to the scholarly life of the Mathematics department,'' May 2019.
\item NSF Graduate Research Fellowship, 2015 -- 2021.
\item Dorothea Klumpke Roberts Prize in Mathematics, Department of Mathematics, UC Berkeley, May 2014.
\item Honors in Mathematics, Department of Mathematics, UC Berkeley, December 2013.
\item Highest Honors, Budapest Semesters in Mathematics, May 2013.
\item Most Outstanding Presentation, Young Mathematicians Conference, Ohio State University, July 2012.
\item Phi Beta Kappa, UC Berkeley Chapter, June 2012.
\end{enumerate}


\section{Teaching}
{\bf University of Utah}
\begin{enumerate}
\item Spring 2024: Math 4400 (introduction to number theory).
\item Fall 2023: Math 5510 (introduction to general topology).
\item August 2023: Co-organizer, new instructor training.
\item Spring 2023: Math 3160 (complex analysis).
\item Fall 2022: Math 1220 (calculus II).
\item Fall 2021: Math 1080 (precalculus).
\end{enumerate}
{\bf University of Michigan}
\begin{enumerate}
\item Summer 2021: Co-instructor, Michigan Math and Science Scholars, course on cryptography.
\item Winter 2021: Math 115 (calculus I).
\item Fall 2020: Course co-coordinator, Math 105 (precalculus).
\item August 2020: Co-organizer, new instructor training.
\item Fall 2018: Math 115 (calculus I).
\item Winter 2018: Grader, Math 632 (algebraic geometry II).
\item Fall 2017: Grader, Math 631 (algebraic geometry I).
\item Winter 2016: Math 105 (precalculus).
\item Fall 2015: Math 105 (precalculus).
\end{enumerate}


\end{resume}
\end{document}
